% Part: turing-machines
% Chapter: machines-computations
% Section: unary-numbers

\documentclass[../../../include/open-logic-section]{subfiles}

\begin{document}

\olfileid[hi]{tur}{mac}{una}
\olsection{संख्याओं का एकाधारी निरूपण}

\begin{explain}
ट्यूरिंग मशीनें अपनी पट्टी पर लिखे प्रतीकों के अनुक्रमों पर काम करती
हैं। ट्यूरिंग मशीन द्वारा प्रयुक्त वर्णमाला के अनुसार ये प्रतीक-अनुक्रम
विभिन्न आगतों और निर्गतों को निरूपित कर सकते हैं। निश्चय ही, प्राकृतिक
संख्याओं के फलनों अर्थात् \emph{अंकगणितीय} फलनों को संगणित करने वाली
ट्यूरिंग मशीनें विशेष रुचि की हैं। धन पूर्णांकों को निरूपित करने की एक
सरल रीति उन्हें एक ही प्रतीक~$\TMstroke$ के अनुक्रमों के रूप में कूटबद्ध
करना है। यदि $n \in \Nat$, तो $\TMstroke^n$ को, $n = 0$ होने पर, रिक्त
अनुक्रम और अन्यथा ठीक $n$ $\TMstroke$ वाला अनुक्रम मानें।
\end{explain}

\begin{defn}[संगणना]
ट्यूरिंग मशीन~$M$, फलन $f\colon \Nat^k \to \Nat$ को \emph{संगणित करती
है}, तभी और केवल तभी जब $M$ आगत
\[
\TMstroke^{n_1} \TMblank \TMstroke^{n_2} \TMblank \dots \TMblank \TMstroke^{n_k}
\]
पर निर्गत $\TMstroke^{f(n_1, \dots, n_k)}$ के साथ रुकती है।
\end{defn}

\begin{prob}
  ट्यूरिंग मशीन~$M$, फलन $f\colon \Nat^k \to \Nat^m$ को कब संगणित
  करती है, इसकी परिभाषा दें।
\end{prob}

\begin{ex}\ollabel{ex:adder}
\emph{योग:}
फलन $f(n,m) = n + m$ को संगणित करने वाली मशीन बनाएँ। इसके लिए ऐसी मशीन
चाहिए जो पट्टी पर $\TMstroke$ के दो खंडों से आरंभ हो, जिनकी लंबाइयाँ
क्रमशः $n$ और~$m$ हों, और $n+m$~$\TMstroke$ वाले एक खंड के साथ रुके। आगत में
~$\TMstroke$ के दोनों खंडों के बीच एक~$\TMblank$ है; अतः एक रीति यह होगी
कि $\TMblank$ वाले खाने पर एक आघात लिख दें और अंतिम~$\TMstroke$ मिटा
दें।
\begin{figure}\[
\begin{tikzpicture}[->,>=stealth',shorten >=1pt,auto,node distance=2.8cm,
                    semithick]
  \tikzstyle{every state}=[fill=none,draw=black,text=black]

  \node[initial,state]         (A)              {$q_0$};
  \node[state]         (B) [right of=A] {$q_1$};
  \node[state]         (C) [right of=B] {$q_2$};

  \path (A) edge node {\TMtrans{\TMblank}{\TMstroke}{\TMstay}} (B)
           edge [loop above] node {\TMtrans{\TMstroke}{\TMstroke}{\TMright}} (B)
        (B) edge [loop above] node {\TMtrans{\TMstroke}{\TMstroke}{\TMright}} (B)
            edge node {\TMtrans{\TMblank}{\TMblank}{\TMleft}} (C)
        (C) edge [loop above] node {\TMtrans{\TMstroke}{\TMblank}{\TMstay}} (C);
\end{tikzpicture}
\]\caption{$f(x,y) = x+y$ संगणित करने वाली मशीन}
\ollabel{fig:adder}
\end{figure}
\end{ex}

\begin{prob}
आगत~$\tuple{3,2}$ के लिए \olref[tur][mac][una]{ex:adder} की मशीन के
विन्यासों का क्रम निकालें। यदि मशीन $0+0$ संगणित करे, तो क्या होता है?
\end{prob}

\begin{explain}
\olref[rep]{ex:doubler} में हमने ऐसी ट्यूरिंग मशीन का उदाहरण दिया था जो
~$\TMstroke$ के अनुक्रम को आगत के रूप में लेती है और पट्टी पर उससे दुगुने
~$\TMstroke$ के अनुक्रम के साथ रुकती है---द्विगुणक मशीन। किंतु निर्गत में
$\TMblank$ भी हैं, जो~$\TMstroke$ के दुगुने खंड के बाईं ओर हैं; इसलिए आपकी
संभावित धारणा के विपरीत, वह वास्तव में फलन $f(x) = 2x$ को संगणित नहीं
करती। इसे ठीक करने की दो रीतियाँ बताएँगे।
\end{explain}

\begin{ex}
\olref{fig:doubler-disc} की मशीन फलन $f(x) =
2x$ को संगणित करती है।
\olref[rep]{ex:doubler} की मशीन की तरह आगत मिटाकर उसके प्रत्येक
$\TMstroke$ के लिए दूर दाईं ओर दो $\TMstroke$ लिखने के स्थान पर यह मशीन
आगत के प्रत्येक~$\TMstroke$ के लिए दाईं ओर एक~$\TMstroke$ जोड़ती है। इसे
लेखा रखना होता है कि आगत कहाँ समाप्त होता है; इसलिए वह आगत और जोड़े गए
आघातों के बीच एक~$\TMblank$ छोड़ती है और सबसे अंत में उसे~$\TMstroke$
से भर देती है। हमें यह भी ``याद'' रखना होता है कि हम आगत में कहाँ हैं;
इसलिए आगत-खंड के एक~$\TMstroke$ को अस्थायी रूप से~$\TMblank$ से बदलते
हैं।
  \begin{figure}
\[
\begin{tikzpicture}[->,>=stealth',shorten >=1pt,auto,node distance=2.8cm,
                    semithick]
  \tikzstyle{every state}=[fill=none,draw=black,text=black]

  \node[initial,state] (0)              {$q_0$};
  \node[state]         (1) [above of=0] {$q_1$};
  \node[state]         (2) [above of=1] {$q_2$};
  \node[state]         (3) [right of=2] {$q_3$};
  \node[state]         (4) [below of=3] {$q_4$};
  \node[state]         (5) [below of=4] {$q_5$};
  \node[state]         (6) [above right of=4] {$q_6$};
  \node[state]         (7) [below of=6] {$q_7$};
  \node[state]         (8) [below of=7] {$q_8$};

  \path (0) edge node {\TMtrans{\TMstroke}{\TMblank}{\TMright}} (1)
%    (0) edge node {\TMtrans{\TMblank}{\TMblank}{\TMstay}} (h)
    (1) edge [loop left] node {\TMtrans{\TMstroke}{\TMstroke}{\TMright}} (1)
      edge node {\TMtrans{\TMblank}{\TMblank}{\TMright}} (2)
    (2) edge [loop above] node {\TMtrans{\TMstroke}{\TMstroke}{\TMright}} (2)
        edge node {\TMtrans{\TMblank}{\TMstroke}{\TMleft}} (3)
    (3) edge [loop above] node {\TMtrans{\TMstroke}{\TMstroke}{\TMleft}} (3)
        edge node[left] {\TMtrans{\TMblank}{\TMblank}{\TMleft}} (4)
    (4) edge        node[left] {\TMtrans{\TMstroke}{\TMstroke}{\TMleft}} (5)
    (5) edge [loop below] node {\TMtrans{\TMstroke}{\TMstroke}{\TMleft}} (5)
        edge              node {\TMtrans{\TMblank}{\TMstroke}{\TMright}} (0)
    (4) edge      node[sloped] {\TMtrans{\TMblank}{\TMstroke}{\TMright}} (6)
    (6) edge              node {\TMtrans{\TMblank}{\TMstroke}{\TMright}} (7)
    (7) edge [loop right] node {\TMtrans{\TMstroke}{\TMstroke}{\TMright}} (7)
        edge              node {\TMtrans{\TMblank}{\TMblank}{\TMleft}} (8)
    (8) edge [loop right] node {\TMtrans{\TMstroke}{\TMblank}{\TMstay}} (8);
    \end{tikzpicture}
\]
\caption{$f(x) = 2x$ संगणित करने वाली मशीन}
\ollabel{fig:doubler-disc}
\end{figure}
\end{ex}

\begin{ex}\ollabel{ex:mover}
$f(x) = 2x$ संगणित करने की दूसरी संभावना यह है कि मूल द्विगुणक मशीन को
रखें, किंतु अंत में ऐसी अवस्थाएँ और निर्देश जोड़ें जो आघातों के दुगुने
खंड को पट्टी के दूर बाएँ सिरे पर ले जाएँ। \olref{fig:mover} की मशीन
केवल यही अंतिम भाग करती है: ~$\TMblank$ के एक खंड और उसके बाद
~$\TMstroke$ के एक खंड वाली पट्टी पर आरंभ करने पर (और शीर्ष को
~$\TMblank$ के खंड में कहीं भी रखने पर), वह $\TMstroke$ को एक-एक करके
मिटाती है और उन्हें पट्टी के आरंभ में लिखती है। काम पूरा होने का पता
लगाने के लिए वह पहले $\TMstroke$ के खंड के अंत को एक~$\TMendtape$ प्रतीक
से चिह्नित करती है, जिसे अंत में मिटा दिया जाता है। हमने अवस्थाओं की
संख्या~$q_6$ से आरंभ की है ताकि इन्हें द्विगुणक मशीन में जोड़ा जा सके।
आपको केवल एक अतिरिक्त निर्देश $\delta(q_0,
\TMblank) = \tuple{q_6,\TMblank,\TMstay}$ चाहिए; अर्थात्~$q_0$ से~$q_6$ तक
~$\TMtrans{\TMblank}{\TMblank}{\TMstay}$ अंकित एक तीर। (एक सूक्ष्म
समस्या है: परिणामी मशीन आगत~$x=0$ के लिए काम नहीं करती। इसे अभ्यास के
लिए छोड़ते हैं।)
\begin{figure}
  \[
  \begin{tikzpicture}[->,>=stealth',shorten >=1pt,auto,node distance=2.8cm,
                      semithick]
    \tikzstyle{every state}=[fill=none,draw=black,text=black]
    \node[initial,state] (6)              {$q_6$};
    \node[state]         (7) [right of=6] {$q_7$};
    \node[state]         (8) [right of=7] {$q_8$};
    \node[state]         (9) [below of=8] {$q_9$};
    \node[state]         (10) [left of=9] {$q_{10}$};
    \node[state]         (11) [left of=10]  {$q_{11}$};
    \node[state]         (12) [below of=11]  {$q_{12}$};
    \node[state]         (13) [right of=12] {$q_{13}$};
    \node[state]         (14) [right of=13] {$q_{14}$};
    \path
    (6)  edge node {\TMtrans{\TMblank}{\TMblank}{\TMright}} (7)
    (7)  edge [loop above] node {\TMtrans{\TMblank}{\TMblank}{\TMright}} (7)
         edge node {\TMtrans{\TMstroke}{\TMstroke}{\TMright}} (8)
    (8)  edge [loop above] node {\TMtrans{\TMstroke}{\TMstroke}{\TMright}} (8)
         edge node {\TMtrans{\TMblank}{\TMendtape}{\TMleft}} (9)
    (9)  edge [loop right] node {\TMtrans{\TMstroke}{\TMstroke}{\TMleft}} (9)
         edge node {\TMtrans{\TMblank}{\TMblank}{\TMright}} (10)
    (10) edge node[above] {\TMtrans{\TMstroke}{\TMblank}{\TMleft}} (11)
    (11) edge [loop above] node {\TMtrans{\TMblank}{\TMblank}{\TMleft}} (11)
         edge node[left] {\begin{tabular}{@{}l@{}}
          \TMtrans{\TMendtape}{\TMendtape}{\TMright}\\
          \TMtrans{\TMstroke}{\TMstroke}{\TMright}
         \end{tabular}} (12)
    (12) edge node[] {\TMtrans{\TMblank}{\TMstroke}{\TMright}} (13)
    (13) edge [loop below] node {\TMtrans{\TMblank}{\TMblank}{\TMright}} (13)
         edge node[sloped] {\TMtrans{\TMstroke}{\TMblank}{\TMleft}} (11)
    (13) edge node {\TMtrans{\TMendtape}{\TMblank}{\TMstay}} (14);
  \end{tikzpicture}
  \]
  \caption{$\TMstroke$ के एक खंड को बाईं ओर ले जाना}
  \ollabel{fig:mover}
\end{figure}
\end{ex}

\begin{prob}
\olref[tur][mac][una]{ex:mover} में हमने ऐसी मशीन बताई जो
\olref[tur][mac][una]{fig:doubler-disc} की द्विगुणक मशीन और
\olref[tur][mac][una]{fig:mover} की स्थानांतरक मशीन का संयोजन है। यदि
इस संयुक्त मशीन को आगत~$x=0$, अर्थात् रिक्त पट्टी पर आरंभ करें, तो क्या
होता है? मशीन को कैसे ठीक करेंगे ताकि इस स्थिति में वह निर्गत~$2x=0$
के साथ रुके? (यह एक अवस्था और एक संक्रमण जोड़कर किया जा सकना चाहिए।)
\end{prob}

\begin{prob}
\emph{घटाव:} ऐसी ट्यूरिंग मशीन की अभिकल्पना करें जो क्रमशः $n$ और~$m$
लंबाई वाली आघातों की दो अरिक्त अक्षर-शृंखलाएँ आगत के रूप में मिलने पर,
जहाँ $n >
m$, फलन $f(n,m) = n - m$ संगणित करे।
\end{prob}

\begin{prob}
\emph{समानता:} निम्न फलन को संगणित करने वाली ट्यूरिंग मशीन की अभिकल्पना
करें:
\[
\fn{equality}(n,m) = 
\begin{cases}
  \text{१} & \text{यदि~$n = m$} \\
  \text{०} & \text{यदि~$n \neq m$}
\end{cases}
\]
जहाँ~$n$ और~$m \in \PosInt$।
\end{prob}

\begin{prob}
फलन $\min(x,y)$ को संगणित करने वाली ट्यूरिंग मशीन की अभिकल्पना करें,
जहाँ $x$ और $y$ ऐसे धन पूर्णांक हैं जिन्हें पट्टी पर $\TMstroke$ की
अक्षर-शृंखलाओं द्वारा निरूपित किया गया है और जो एक $\TMblank$ से अलग हैं। मशीन की
वर्णमाला में अतिरिक्त प्रतीकों का उपयोग कर सकते हैं।

फलन $\min$ अपने तर्कों में सबसे छोटा मान चुनता है; अतः
$\min(3,5)=3$, $\min(20,16)=16$, और $\min(4,4)=4$, इत्यादि।
\end{prob}

\begin{defn}
  ट्यूरिंग मशीन~$M$, आंशिक फलन $f\colon \Nat^k
  \pto \Nat$ को संगणित
  करती है, तभी और केवल तभी जब
  \begin{enumerate}
    \item $M$ आगत $\TMstroke^{n_1}\concat\TMblank\concat
    \dots \concat\TMblank\concat\TMstroke^{n_k}$ पर निर्गत
    $\TMstroke^{m}$ के साथ रुकती है, यदि $f(n_1, \dots, n_k) = m$।
    \item $M$ बिल्कुल नहीं रुकती, अथवा ऐसे निर्गत के साथ रुकती है
    जो~$\TMstroke$ का एकल खंड नहीं है, यदि $f(n_1, \dots, n_k)$
    अपरिभाषित है।
  \end{enumerate}
\end{defn}

\end{document}
